\scanpage{15}
\[L_{DE}(t)=L_E\left(p^{-(n+\rho)}t\right)\]
d'où, portant dans (2) et remplaçant $t$ par $t^{-1}$ :
\[
(3)\qquad\boxed{L_E\left(\frac1{p^{n+\rho}t}\right)=\varepsilon(E)\left(-p^{\frac{n+\rho}{2}}t\right)^{\chi(E)}L_E(t)}
\]
Posant
\[\boxed{\unsure{\xi|\zeta}_E(t)=t^{\chi(E)/2}L_E(t)}\]
il vient l'équation fonctionnelle
\[\boxed{\unsure{\xi|\zeta}\left(\frac1{p^{n+\rho}t}\right)=\varepsilon'(E)\unsure{\xi|\zeta}(t)}\]
\[
\begin{aligned}
\varepsilon'(E)&=\pm1,\qquad\varepsilon'(E)=(-1)^{\chi(E)}\prod_i\operatorname{signe}\det f_{H^i(\bar X,\bar F)}\\
&\hspace{7em}=(-1)^{\chi(E)}\operatorname{signe}\det f_{H^*(\bar X,\bar F)}
\end{aligned}
\]
\deleted{Si on écrit l'équation fonctionnelle avec $\varepsilon(E)$ au lieu de}

Notons d'ailleurs que $H^i(\bar X,\bar E)$ et $H^{2n-i}(\bar X,\check{\bar E})$ \inserted{$\simeq H^{2n-i}(\bar X,\bar E(\rho))=H^{2n-i}(\bar X,\bar E)(\rho)$} étant ``duaux à valeurs dans $\mathbb Q_\ell(n)$'', et l'opération de $f$ sur \inserted{$\mathbb Q_\ell(n)$} \unsure{ce dernier positive},
\[\varepsilon_i(E)=\varepsilon_{2n-i}(E),\]
et par suite on trouve que
\[\boxed{\varepsilon(E)=\operatorname{Signe}\det f_{H^n(\bar X,\bar E)}}.\]
\unsure{Pour la raison analogue}, $(-1)^{\chi(E)}=(-1)^{b_n}$. Donc
\[\varepsilon'(E)=(-1)^{b_n(E)}\operatorname{signe}\det f_{H^n(\bar X,\bar E)}.\]
\mspage{14}
